\scanpage{109}
\deleted{\uline{Corollaire 1}\quad Supposons que dans $X$ et $Y$, \deleted{\unsure{tt pt}} \deleted{\unsure{admettent} \unsure{une|un} \unsure{système|application} \unsure{fondamental} de \unsure{vois.}}\inserted{base d'ouverts \inserted{$X_i$ resp. $Y_j$} rétrocompacts quasi-compacts}. Alors l'hyp préliminaire de la prop. équivaut aussi à dire que Pour \unsure{$\forall$} ouvert \deleted{\unsure{quasi}} $Y_j$, $f^{-1}(Y_j)$ est rétrocompact dans $X$, ou encore que pour \unsure{$\forall$} \deleted{ouvert} ouvert $Y'$ \deleted{quasi cpt}\inserted{\deleted{\unsure{E rétrocompact}} rétrocompact} de $Y$, $X'=f^{-1}(\unsure{Y'|U'})$ est \deleted{q. cpt} rétrocompact dans $X$, ou encore que les $X_i\cap f^{-1}(Y_j)$ sont q. cpts.}

[\uline{Prop.} Soit $X$ un espace topologique, ayant une base d'ouverts $U_i$ qui sont quasi-compacts. Alors pour \unsure{$\forall$} \deleted{partie $E$ de $X$, il revient} application continue $f:X'\unsure{-|\longrightarrow}X$, il revient au même de dire que $f$ est q. cpt, ou que les $f^{-1}(U_i)$ sont q. cpts.

Donc si $X'\subset X$ sous-espace top. $X'$ est rétrocompact dans $X$ ssi les $X'\cap U_i$ sont q. cpts.

\uline{Cor.} \deleted{Soit \unsure{Supposons} de plus}\inserted{Pour que} les $U_i$ rétrocompacts dans $X$, \deleted{Alors}\inserted{il f. et s. que} \uline{\unsure{$\forall$}} partie quasi-compacte ouverte de $X$ soit rétrocompacte dans $X$.

\margin{\deleted{\uline{Cor} b) $f:X'\unsure{-|\longrightarrow}X$ q. cpt \unsure{dans} $X$ \illegible{2}}}

\uline{Cor.} S'il en est ainsi, pour une \deleted{\unsure{ap}} applic. continue $f:X'\unsure{-|\longrightarrow}X$, il revient au même de \unsure{dire} \unsure{conditions}

(a) pour \unsure{$\forall$} $U\subset X$ \inserted{ouvert} qui est rétrocompact, $f^{-1}(U)$ est rétrocompact,

(b) Pour \unsure{$\forall$} $U_i$, $f^{-1}(U_i)$ est rétrocompact dans $X'$,

(c) [Si $X'$ a une base d'ouverts quasi-compacts $U'_j$] les $f^{-1}(U_i)\cap U'_j$ sont \deleted{rétro} qua cpts.]

\margin{$U\subset X$ ouv. rétrocompact, si $U'\subset X'$ ouv. q. cpt}
\margin{$b\Rightarrow a$, En effet, $f^{-1}(U)\cap\unsure{V'|U'}\subset f^{-1}(f(U'))$, \deleted{$=f^{-1}(U)\cap\unsure{f^{-1}}$}, et $f(U')$ étant q. cpt, est contenu dans une réunion finie de $U_i$ $(i\in J)$, donc
\[
f^{-1}(U)\cap V'=f^{-1}\left(\bigcup_{i\in J}U\cap U_i\right)\cap V'.
\]
Or chaque $U\cap U_i$ est réunion finie \unsure{d'uns} $U_j$}

\margin{\figure{\begin{tikzpicture}[x=1cm,y=1cm]
\draw (0,0)--(0.15,-2.1);
\draw (0.3,-4.5)..controls(0.3,-5.1)and(0.7,-5.1)..(1.5,-5.2)..controls(2.7,-5.25)and(3.1,-5.5)..(3.35,-5.9);
\end{tikzpicture}}}
\margin{108}
